% Compile twice with XeLaTeX; requires xeCJK and Noto CJK fonts.
\documentclass[10pt,a4paper]{article}
\usepackage[a4paper,left=18mm,right=18mm,top=16mm,bottom=17mm]{geometry}
\usepackage{fontspec}
\defaultfontfeatures{Ligatures=TeX}
\IfFontExistsTF{Lato}{\setmainfont{Lato}}{\setmainfont{TeX Gyre Heros}}
\usepackage{xeCJK}
\IfFontExistsTF{Noto Sans CJK SC}{\setCJKmainfont{Noto Sans CJK SC}}{\setCJKmainfont{FandolHei-Regular}[BoldFont=FandolHei-Bold]}
\IfFontExistsTF{Noto Serif CJK SC}{\newCJKfontfamily\displayfont{Noto Serif CJK SC}}{\newCJKfontfamily\displayfont{FandolSong-Regular}[BoldFont=FandolSong-Bold]}
\usepackage{xcolor}
\definecolor{ink}{HTML}{203238}
\definecolor{accent}{HTML}{316C6D}
\definecolor{muted}{HTML}{647175}
\definecolor{rulecolor}{HTML}{C7D7D5}
\usepackage[protrusion=false]{microtype}
\usepackage{enumitem}
\usepackage{fancyhdr}
\usepackage{lastpage}
\usepackage{hyperref}
\hypersetup{colorlinks=true,urlcolor=accent,linkcolor=accent,
  pdftitle={熊天宇 - 学术简历},pdfauthor={Tianyu Xiong},
  pdfsubject={三维计算机视觉、学术论文与科研经历},pdflang={zh-CN}}
\urlstyle{same}
\setlength{\parindent}{0pt}
\setlength{\parskip}{0pt}
\setlength{\emergencystretch}{2em}
\renewcommand{\baselinestretch}{1.16}
\color{ink}
\raggedright
\pagestyle{fancy}
\fancyhf{}
\renewcommand{\headrulewidth}{0pt}
\renewcommand{\footrulewidth}{0pt}
\fancyfoot[L]{\fontsize{7.6}{9}\selectfont\color{muted}熊天宇\quad /\quad 学术简历}
\fancyfoot[R]{\fontsize{7.6}{9}\selectfont\color{muted}2026 年 10 月\quad\thepage\,/\,\pageref*{LastPage}}
\setlength{\footskip}{24pt}
\newcommand{\cvsection}[1]{\par\vspace{10pt}{\fontsize{11}{14}\selectfont\bfseries\color{accent}#1}\par\vspace{4pt}{\color{rulecolor}\hrule height .55pt}\vspace{7pt}}
\newcommand{\entry}[3]{\textbf{#1}\hfill{\small\color{muted}#2}\par{\small\color{muted}#3}\par\vspace{4pt}}
\newcommand{\paper}[4]{\textbf{#1}\par\vspace{2pt}{\small #2}\par\vspace{1pt}{\small\color{muted}#3}\quad{\small #4}\par\vspace{9pt}}
\newcommand{\project}[2]{\vspace{4pt}\textbf{#1}\quad{\small\color{muted}#2}\par\vspace{4pt}}
\newcommand{\detail}[1]{{\small #1}\par}
\setlist[itemize]{leftmargin=13pt,itemsep=5pt,topsep=5pt,parsep=0pt,label={\color{accent}\textbullet}}
\begin{document}

{\displayfont\fontsize{32}{40}\selectfont 熊天宇}\quad{\fontsize{18}{22}\selectfont\color{muted}Tianyu Xiong}\par
\vspace{6pt}
{\fontsize{10.5}{14}\selectfont\bfseries\color{accent}博士研究生\quad /\quad 三维计算机视觉}\par
\vspace{9pt}
{\small
\href{mailto:tianyuxiong1997@gmail.com}{tianyuxiong1997@gmail.com}\quad\textbar\quad
\href{https://github.com/Six-Bit-TX}{github.com/Six-Bit-TX}\par
\vspace{3pt}
\href{https://six-bit-tx.github.io/research/zh.html}{six-bit-tx.github.io/research/zh.html}}

\cvsection{个人简介}
西北工业大学计算机科学与技术专业博士研究生。论文涉及三维重建、高斯泼溅、水印与机器学习。

\cvsection{教育背景}
\entry{西北工业大学}{2024 年至今}{计算机科学与技术专业，博士研究生\hfill 中国，西安}
\detail{导师：张艳宁（Yanning Zhang）、Yangjia Qi。}
\vspace{8pt}
\entry{多伦多大学}{2018-2021 年}{荣誉理学学士\hfill 圣乔治校区，加拿大，多伦多}
\detail{专业方向：Computer Science Specialist。}

\cvsection{学术论文}
\paper{Queryable Gaussian View Index: Text Retrieval and Selective Camera Localization}
{\textbf{Tianyu Xiong}, Rui Li, Jiaqi Yang.}
{稿件，2026 年。已提交 arXiv。}
{}

\paper{SalientGS: Unified SfM-to-3DGS with Importance-Guided MCMC Gaussian Allocation}
{\textbf{Tianyu Xiong}, Rui Li, Suning Ge, Jiaqi Yang.}
{ACM Multimedia 2026，已录用。}
{\href{https://arxiv.org/abs/2607.11285}{arXiv:2607.11285}}

\paper{CAGE-LIVO: Cooperative Air-Ground Gaussian-Enhanced LiDAR-Inertial-Visual Odometry}
{\textbf{Tianyu Xiong} et al.（第一作者）}
{Chinese Journal of Aeronautics，已录用。}
{}

\paper{VoMarkSplat: Robust watermarking for 3D Gaussian splatting with patch and multi-convolutional voting}
{\textbf{Tianyu Xiong}, Linjie Li, Rui Li, Jiaqi Yang, Yanning Zhang.}
{Pattern Recognition Letters, 203:8-14，2026 年 5 月。}
{\href{https://doi.org/10.1016/j.patrec.2026.02.024}{DOI}}

\paper{GloSplat: Joint Pose-Appearance Optimization for Faster and More Accurate 3D Reconstruction}
{\textbf{Tianyu Xiong}, Rui Li, Linjie Li, Jiaqi Yang.}
{预印本，2026 年。}
{\href{https://arxiv.org/abs/2603.04847}{arXiv:2603.04847}}

\paper{Large scale air pollution prediction with deep convolutional networks}
{Gao Huang, Chunjiang Ge, \textbf{Tianyu Xiong}, Shiji Song, Le Yang, Baoxian Liu, Wenjun Yin, Cheng Wu.}
{Science China Information Sciences, 64(9)，文章号 192107，2021 年。}
{\href{https://doi.org/10.1007/s11432-020-2951-1}{DOI}}

\paper{The effect of similarity on the evolution of fairness in the ultimatum game}
{Yakun Zhao, \textbf{Tianyu Xiong}, Lei Zheng, Yumeng Li, Xiaojie Chen.}
{Chaos, Solitons \& Fractals, 131，文章号 109494，2020 年 2 月。}
{\href{https://doi.org/10.1016/j.chaos.2019.109494}{DOI}}

\cvsection{技术方法与语言能力}
\detail{\textbf{技术方法：}高斯泼溅、可微渲染、多视几何、神经网络学习与基于 CLIP 的视觉语言模型。}
\detail{\textbf{英语：}母语水平。}

\newpage
{\displayfont\fontsize{23}{31}\selectfont 研究与经历}\par
\vspace{4pt}
{\small\color{muted}熊天宇\quad /\quad 三维计算机视觉与机器学习}

\renewcommand{\detail}[1]{{\fontsize{10.5}{17}\selectfont #1}\par}
\renewcommand{\cvsection}[1]{\par\vspace{20pt}{\fontsize{11}{14}\selectfont\bfseries\color{accent}#1}\par\vspace{5pt}{\color{rulecolor}\hrule height .55pt}\vspace{11pt}}
\cvsection{已完成研究选介}
\project{SalientGS}{快速重建与容量分配}
\detail{构建统一重建流程，对相机位姿与高斯外观进行联合优化。通过重要性引导的 MCMC 分配，将固定高斯预算优先用于持续欠拟合的区域。}
\vspace{13pt}
\project{VoMarkSplat}{三维高斯资产水印}
\detail{通过空间分块投票、多卷积投票及两阶段课程训练，实现高斯场景资产的鲁棒水印恢复。}

\cvsection{早期科研经历}
\entry{香港理工大学}{2022 年 5 月至 2024 年 5 月}{电子计算学系\hfill 中国，香港}
\detail{导师：Lei Zhang；与 Hongwei Yong 合作开展研究项目。}
\begin{itemize}
\item 构建基于 CLIP 的图像裁剪原型，支持内容与美学指令引导。
\item 研究卷积神经网络反向传播中的特征子采样，以及基于直接反馈对齐（Direct Feedback Alignment）的可并行反馈机制。
\end{itemize}
\vspace{10pt}
\entry{清华大学}{2019 年 5 月至 8 月}{自动化系\hfill 中国，北京}
\detail{导师：Gao Huang。参与基于深度卷积网络的时空空气污染预测研究，相关论文于 2021 年发表于 Science China Information Sciences。}
\vspace{14pt}
\entry{北京航空航天大学}{2018 年 5 月至 8 月}{电子信息工程学院\hfill 中国，北京}
\detail{导师：Xiantong Zhen。实现 CliqueNet，并用于 MURA 放射影像分类。}

\cvsection{创业经历}
\entry{Halfman Tech / 联合创始人}{2024 年 5 月至 9 月}{中国，北京}
\detail{联合创办早期科技企业；公司于 2024 年 9 月停止运营。}

\end{document}
